\section{Ethics and Reproducibility}

\paragraph{Data and actions.}
Default tasks use synthetic mock data. The simulated email tool creates drafts only; it never sends real email. The booking tool is a stub and never makes real reservations. The benchmark should not include private personal data, API keys, proprietary documents, or live credentials.

\paragraph{Responsible release.}
\benchmark can be used to diagnose robustness failures, but it can also be gamed if treated as a static leaderboard. Releases should be versioned, include held-out components when possible, and distinguish oracle sanity checks from realistic agents. Leaderboard exports should exclude \texttt{scripted\_oracle\_agent} from model rankings by default; see the reviewer attack matrix (\texttt{reviews/reviewer\_attack\_response\_matrix.md}, attack~17).

\paragraph{Human subjects.}
If human validation is used, annotators should receive clear instructions, fair compensation, consent information, and a process for reporting ambiguous or uncomfortable examples. The paper must report the validation protocol and agreement statistics.

\paragraph{Compensation placeholder.}
The repository includes protocol placeholders for annotator qualifications, expected time per item, compensation rate, consent language, and adjudication responsibilities. These must be completed before collecting labels used in paper claims.

\paragraph{Transparency.}
Every reported result should include task version, run config, random seed, model/API version, timestamp, scorer version, and artifact path. Current reproducibility conventions are documented in \texttt{docs/REPRODUCIBILITY.md}. Empirical claim requirements and forbidden wording before experiments complete are tracked in \texttt{docs/claim\_ledger.json}, \texttt{paper/EVIDENCE\_GAP\_MAP.md}, and \texttt{docs/DO\_NOT\_OVERCLAIM.md}. The smoke-test and mock-pipeline reproducibility claim is tracked by \claimref{C9}; broader scientific claims remain planned until the main experiments are complete.

\paragraph{Evidence-level enforcement.}
Runs record \texttt{evidence\_scope}, \texttt{scientific\_evidence}, and \texttt{paid\_calls\_made} in metadata. Validators (\texttt{check\_evidence\_safety.py}, \texttt{check\_claim\_ledger.py}) block incomplete or mock/stub runs from supporting paper claims. Paid API calls are disabled by default (\texttt{allow\_paid\_calls: false}).

\paragraph{Release and environment capture.}
Release packaging includes a versioned manifest (\texttt{build-release-manifest}), a reproducibility bundle plan (\texttt{plan-repro-bundle}), and environment capture without secrets (\texttt{capture-env}). These artifacts support artifact evaluation and future public releases; they do not imply that main experimental results are complete.

\paragraph{Pre-experiment gates.}
Before provider or main experiments, the pre-experiment freeze checklist (\texttt{experiments/PRE\_EXPERIMENT\_FREEZE\_CHECKLIST.md}) and main experiment gate (\texttt{experiments/MAIN\_EXPERIMENT\_GATE.md}) must pass. This draft reports framework readiness, not gate completion for main results.
